\paragraph{AQ-136-SEMILOCAL: arithmetic path gluing for the ordinary localization.}
Let $P=\{p_1,\ldots,p_s\}$ be a finite nonempty set of rational
primes, and put
\[
 R_P=\mathbb Z[\ell^{-1}:\ell\text{ prime},\ \ell\notin P].
\]
For every commutative $R_P$-algebra $k$, the ordinary groupoid
obtained from $\mathcal A_k$ by inverting every morphism is
equivalent to the terminal groupoid. Together with the rational
case, this proves the assertion for every commutative ring in
which all but finitely many rational primes are units.

This is a theorem about the original localization question.
It does not assert a strict common admissible source for two
given theories over $R_P$ when $|P|>1$. Its proof constructs
common sources after specified natural arrows to arithmetic
path replacements. The integral ring $\mathbb Z$, and rings
with infinitely many nonunit rational primes, are not included.
The zero ring is harmless: its unital algebra category has
only the zero algebra. Assume henceforth that $k\ne0$.

\paragraph{1. An exact arithmetic sequence for arbitrary modules.}
Set
\[
 k_i=k\otimes_{R_P}\mathbb Z_{(p_i)},\qquad
 K=k\otimes_{R_P}\mathbb Q.
\]
These are localizations of $k$: to obtain $k_i$, invert the
product of the primes in $P\setminus\{p_i\}$, and to obtain
$K$, invert the product of all primes in $P$.
For every $k$-module $M$, write $M_i=M\otimes_k k_i$ and
$M_0=M\otimes_k K$. There is an exact sequence
\begin{equation}
 \begin{aligned}
 0\longrightarrow M\longrightarrow
       \left(\bigoplus_{i=1}^sM_i\right)\oplus M_0
   &\xrightarrow{\delta} M_0^{\oplus s}\longrightarrow0,\\
 \delta((m_i),m_0)&=(m_i-m_0)_i.
 \end{aligned}                                           \tag{136.S1}
\end{equation}

Here is a direct proof, including the base-change assertion.
For $M=R_P$, the kernel statement is
$\bigcap_i\mathbb Z_{(p_i)}=R_P$ inside $\mathbb Q$.
For surjectivity, given finitely many $a_i\in\mathbb Q$,
choose a common denominator
$D=d\prod_i p_i^{e_i}$ with $d$ prime to every $p_i$.
Write $a_i=b_i/D$. The Chinese remainder theorem gives an
integer $b$ such that $b+b_i$ is divisible by $p_i^{e_i}$
for every $i$. With $a_0=b/D$ we have
$a_i+a_0\in\mathbb Z_{(p_i)}$, as required.
The rightmost module $\mathbb Q^{\oplus s}$ is flat over
$R_P$, so tensoring this short exact sequence with any
$R_P$-module remains exact. Apply this to the underlying
$R_P$-module of $M$. This proves (136.S1) for arbitrary
$M$, with no flatness assumption on $k$ or on $M$.

All localizations used below are flat. Thus vertical and
normalized horizontal homology commute with them, and an
admissible theory localizes to an admissible theory with
the same weight cutoffs. Normalization commutes with these
base changes, either by exactness or by its splitting
idempotents. The coefficient algebra localizes as
$\Gamma_k(T)\otimes_k k_i=\Gamma_{k_i}(T)$.

\paragraph{2. A rational tensor lemma with a uniform weight bound.}
Let $K$ be any $\mathbb Q$-algebra, possibly zero, and let
$N_1,\ldots,N_s$ be positive integers. Over $K$, let $B_{N_i}$
be the weighted simplex theory whose zero-form coordinates
have weight $N_i$ and whose coefficient relation is
$\sum_jx_{i,j}=T^{[N_i]}$. It is the rationalization of
the basic source in AQ-136-PLOCAL. Set
\[
 F=\bigotimes_{i=1}^s{}_{\Gamma_K(T)}B_{N_i}.
\]
Then $F$ is admissible. Its vertical cohomology is exactly
$\Gamma_K(T)$ in degree zero in every weight; for fixed
form degree $q$, its horizontal homology vanishes in all
weights
\[
                    w\ge q\max_iN_i+\sum_iN_i.           \tag{136.S2}
\]

We give the module calculation rather than infer this claim
from a vertical K\"unneth argument. Identify
$\Gamma_K(T)$ with $S_K=K[t]$, by $T^{[a]}=t^a/a!$.
Rescaling the coordinates by $N_i!$ writes the $i$th
coefficient relation as $\sum_jx_{i,j}=t^{N_i}$.
Before this base change, use the coefficient algebra
$K[U_i]$ and relation $U_i=\sum_jx_{i,j}$, where $U_i$
and the coordinates have weight $N_i$. Its form-$q_i$
row is
\[
 \bigoplus_{m\ge0}\operatorname{Sym}^m_K(L)
                         \otimes_K\Lambda^{q_i}_K(Q),
 \qquad Q=L/K\!\left(\sum_jx_j\right),                 \tag{136.S3}
\]
with weight $N_i(m+q_i)$ in the $m$th summand.
The positive-$m$ summands are simplicially contractible:
apply $\operatorname{Sym}^m$ to the explicit linear
contraction of $L$, and tensor the resulting simplicial
homotopy with the identity on $\Lambda^{q_i}Q$.
Hence the augmentation of this row onto
$\Lambda^{q_i}Q$, on which $U_i$ acts as zero, is a
horizontal quasi-isomorphism. Each simplex of the row
is a free $K[U_i]$-module: eliminate $x_{i,0}$ using
$U_i$, and use the exterior basis on the remaining
differentials.

For a fixed tuple $(q_1,\ldots,q_s)$ with sum $q$, take
the tensor of the rows (136.S3), using the independent
coefficient ring $A=K[U_1,\ldots,U_s]$.
The same contraction, on any factor with positive $m$,
shows that it maps by a horizontal quasi-isomorphism to
\[
                V_* =\bigotimes_i\Lambda^{q_i}Q_*.
\]
All of $V_*$ has the single internal weight
$b=\sum_iq_iN_i$, and every normalized term of $V_*$
is projective over $K$, being a direct summand of a
finite free module. The normalized source is a
nonnegative complex of projective $A$-modules. It
therefore computes the derived base change
$S_K\otimes_A^{\mathbf L}N(V_*)$.
Here ordinary base change of the normalized source is
the normalized row after base change; the split
normalization idempotents justify this even though
$A\to S_K$ need not be flat.

Resolve $K=A/(U_1,\ldots,U_s)$ by the Koszul complex
on the variables $U_i$. It is exact over arbitrary
$K$: successive multiplication by a polynomial variable
is injective. Tensoring this resolution with the
nonnegative projective $K$-complex $N(V_*)$ is a
projective $A$-resolution of that complex, by the
first-quadrant filtration. Thus the required derived
base change is computed by
\[
           K(t^{N_1},\ldots,t^{N_s};S_K)
                               \otimes_K N(V_*).         \tag{136.S4}
\]
Reorder the indices so $N_1=\min_iN_i$, and let $\xi_i$
be the Koszul generator of weight $N_i$ and homological
degree one. The invertible basis change
\[
        \eta_i=\xi_i-t^{N_i-N_1}\xi_1\quad(i>1)
\]
has $d\eta_i=0$. Consequently (136.S4) is
quasi-isomorphic to
\[
  S_K/(t^{N_1})\otimes_K
       \Lambda_K(\eta_2,\ldots,\eta_s)\otimes_K N(V_*).
                                                               \tag{136.S5}
\]
The passage to the quotient is valid after tensoring:
$N(V_*)$ is bounded below with projective $K$-terms,
and the exterior factors are free. Every term in
(136.S5) has weight at most
\[
       b+(N_1-1)+\sum_{i>1}N_i
                    \le q\max_iN_i+\sum_iN_i-1.
\]
There are finitely many tuples $(q_i)$ for a fixed $q$,
so this proves (136.S2). The tensor-to-normalized
comparisons used here are the flat Eilenberg--Zilber
comparison of Stacks, Remark 21.39.11, Tag \texttt{08QD};
equivalently one can apply its first-quadrant total
complex form at each finite tensor step.

Finally, simplexwise, eliminate one zero-form coordinate
per factor using $t^{N_i}$. The vertical complex is
$S_K$ tensored with finitely many polynomial de Rham
disks. Integration over the $\mathbb Q$-algebra $K$
contracts each disk onto its constants, preserving
weight. This proves the stated vertical comparison.
Thus $F$ is admissible, with a bound independent of
the simplex dimension. This lemma includes $s=1$.

\paragraph{3. Arithmetic path replacements.}
For a differential $K$-algebra $D$, put
\[
 \operatorname{Path}_{\rm dR}(D)=D\otimes_K K[z,dz],
 \quad |z|=0,\quad |dz|=1,\quad\operatorname{wt}(z)
                  =\operatorname{wt}(dz)=0.
\]
Evaluation at $z=0,1$, $dz=0$ gives $e_0,e_1$.
Integration is a natural contraction of this path
complex onto $D$, so both evaluations induce the
same isomorphism in vertical cohomology. This
construction is used only over the $\mathbb Q$-algebra
$K$. Its positive-weight divided powers are uniquely
specified by $\gamma_m(a)=a^m/m!$ on even elements.

For $C\in\mathcal A_k$, write $C_i=C\otimes_k k_i$
and $C_0=C\otimes_k K$. Define $\mathcal E(C)$,
in every simplex degree, to consist of tuples
\[
       ((c_i)_{i=1}^s,c_0,(\omega_i)_{i=1}^s)
\]
with $c_i\in C_i$, $c_0\in C_0$, and
$\omega_i\in\operatorname{Path}_{\rm dR}(C_0)$,
subject to
\[
                       e_0\omega_i=c_i|_K,
             \qquad    e_1\omega_i=c_0.                   \tag{136.S6}
\]
This is a limit of strict bigraded PD differential
algebras, so it inherits all the required operations.
It is simplicial and functorial in $C$. The coefficient
map takes a coefficient to its localizations and
constant paths. The same prescription gives a natural
morphism
\[
                            \eta_C:C\longrightarrow\mathcal E(C).
                                                               \tag{136.S7}
\]

We verify admissibility explicitly. Forgetting
multiplication gives a short exact sequence of
simplicial cochain modules
\[
0\longrightarrow\mathcal E(C)\longrightarrow D(C)
 \longrightarrow C_0^{\oplus 2s}\longrightarrow0,          \tag{136.S8}
\]
where
$D(C)=(\bigoplus_i C_i)\oplus C_0
\oplus\operatorname{Path}_{\rm dR}(C_0)^{\oplus s}$,
and the last map is the endpoint difference in
(136.S6). It is surjective in every cochain and
simplex degree: $(1-z)a+zb$ gives prescribed
endpoints $(a,b)$. This section is linear and
simplicially natural; it need not be a cochain map.

Fix a cochain cutoff $d$. At sufficiently large weight,
uniformly in simplex degree, every localized vertical
complex has coefficient cohomology in degree zero
and zero cohomology in degrees $1$ through $d$.
The same is true of the rational path complexes.
On degree-zero cohomology, the cokernel of the map
in (136.S8), after eliminating the path constants,
is exactly the cokernel of (136.S1), and its kernel
is $S_k^w$. Thus the long exact vertical cohomology
sequence gives $H^0=S_k^w$ and $H^j=0$ for
$1\le j\le d$. All bounds are independent of
simplex degree. For fixed form degree $q$, the
path row is a direct sum of copies of the two rows
$C_0^{w,q}$ and $C_0^{w,q-1}$. Polynomial coefficients
are free over $K$, so these rows are horizontally
acyclic at large weight. The other rows in (136.S8)
are likewise acyclic, and its horizontal long exact
sequence proves the same for $\mathcal E(C)$.
This proves $\mathcal E(C)\in\mathcal A_k$.

We will also use its compatibility with the simplicial
cotensor $P_s$ of AQ-136-COTENSOR:
\[
                     \mathcal E(P_s C)\simeq P_s\mathcal E(C).
                                                               \tag{136.S9}
\]
For each $n$, the simplicial set $\Delta[n]\times\Delta[1]$
has finitely many nondegenerate simplices. Sections of
a simplicial module on it are therefore computed by
a finite limit of values on simplices. Flat localization
commutes with these limits. Tensoring with $K[z]$ also
commutes with them, since it is flat, and a fixed path
cochain degree is a sum of only the two displayed
form degrees. Limits commute with the finite limits
in (136.S6). These natural identifications preserve
products, the differential, divided powers and all
endpoint maps, proving (136.S9).

\paragraph{4. One local source and homotopies for an entire finite diagram.}
Let a finite directed graph be labelled by objects
$C_v\in\mathcal A_k$ and actual morphisms
$f_e:C_v\to C_{v'}$ on its edges. For each $i$, choose
a power $N_i=p_i^{a_i}$ so large that every zero-form
row $(C_v)_i^{N_i,0}$ is horizontally acyclic.
There is one such power for all vertices, since they
are finite in number. The p-local construction gives
the basic admissible source
$A_i=B_{N_i}(k_i)$ and maps
\[
                         a_{i,v}:A_i\longrightarrow(C_v)_i
                                                               \tag{136.S10}
\]
for every vertex.

For each edge, the two maps $f_ea_{i,v}$ and
$a_{i,v'}$ from this \emph{same} $A_i$ to $(C_{v'})_i$
have a coefficient-preserving simplicial PD homotopy
\[
                  h_{i,e}:A_i\longrightarrow P_s((C_{v'})_i).
                                                               \tag{136.S11}
\]
Here is the simultaneous-choice justification. On the
weight-$N_i$ coordinate module $L$, the difference of
the two coordinate maps kills the constant submodule
$k_i\sum x_j$, because both maps take it to $T^{[N_i]}$.
It factors through $Q=L/k_i\sum x_j$, whose normalized
complex is $k_i[1]$. Horizontal $H_1$ of the target
zero-form row vanishes. The resulting degree-one
cycle is therefore a boundary, giving a chain
nullhomotopy, and Dold--Kan gives a simplicial linear
homotopy. Adding the first coordinate map makes it
relative to the coefficient submodule. The free PD
coordinate presentation extends it, with its
differential, to (136.S11). This is precisely the
construction checked in AQ-136-PLOCAL-LOCALIZATION.
The same row cutoff works for every edge; choices for
the finitely many homotopies may be made separately
after the single common $A_i$ and all maps (136.S10)
have been fixed. The resulting data supply every
edge of the finite graph at once.

\paragraph{5. Gluing that finite diagram and killing every finite loop.}
Form the rational admissible theory
\[
                    F=\bigotimes_{i=1}^s{}_{S_K}(A_i\otimes_{k_i}K),
\]
which is covered by Step 2, and let $j_i:(A_i)_K\to F$
be the canonical inclusions. Define $A$ by the same
path-gluing formula as (136.S6), using local objects
$A_i$, central object $F$, and endpoint maps $j_i$:
\[
 A=\left\{((a_i),a_0,(\omega_i)):
     \begin{array}{l}
       e_0\omega_i=j_i(a_i|_K),\\
       e_1\omega_i=a_0\quad(1\le i\le s)
     \end{array}\right\}.                                 \tag{136.S12}
\]
Here $a_i\in A_i$, $a_0\in F$, and
$\omega_i\in\operatorname{Path}_{\rm dR}(F)$.
The verification in Step 3 applies verbatim to this
diagram of different theories. In detail, its endpoint
difference map is still degreewise surjective; all
local and central vertical cohomology groups are
exactly their coefficient rings, and the induced
maps are the coefficient localizations. Sequence
(136.S1) consequently gives vertical cohomology
exactly $S_k$ in every weight. Its horizontal rows
are eventually acyclic by the finitely many p-local
source cutoffs and the uniform rational bound (136.S2).
Hence $A\in\mathcal A_k$.

For each vertex, multiplication of the rationalized
maps (136.S10) defines a strict map
\[
        b_v:F\longrightarrow(C_v)_0,\qquad
        b_v(u_1\otimes\cdots\otimes u_s)
                   =\prod_i a_{i,v}(u_i).
\]
It is well defined over $S_K$, since all component
maps agree on that coefficient algebra. It preserves
the differential and the gradings; over the rational
algebra it also preserves the uniquely specified
divided powers. Its restriction along $j_i$ is
$(a_{i,v})_K$. Applying it to the central value and
the paths in (136.S12), with (136.S10) on the local
values, gives a strict map
\[
                         \alpha_v:A\longrightarrow\mathcal E(C_v).
                                                               \tag{136.S13}
\]

Likewise multiply the rationalized maps (136.S11)
to obtain $b_e:F\to P_s((C_{v'})_0)$. Its endpoints
are $f_eb_v$ and $b_{v'}$: this follows from
multiplicativity of $f_e$ and of the endpoint maps.
On local values use $h_{i,e}$, on the central value
use $b_e$, and on rational paths use its differential
path extension. This gives
\[
 A\longrightarrow\mathcal E(P_s C_{v'})
                         \simeq P_s\mathcal E(C_{v'}),
\]
whose endpoints are $\mathcal E(f_e)\alpha_v$ and
$\alpha_{v'}$. Thus the single source $A$ supplies
homotopies for every edge of the entire finite graph,
with all its vertex maps fixed.

In the all-arrow localization, the two evaluations
$P_sD\rightrightarrows D$ are equal: they are both
inverse to the constant-path arrow $D\to P_sD$.
All three objects and arrows are admissible by
AQ-136-COTENSOR. Consequently
\[
                   \mathcal E(f_e)\alpha_v=\alpha_{v'}
\]
in that localization. Naturality of (136.S7) now
shows that, with $\beta_v=\eta_{C_v}^{-1}\alpha_v$,
\[
                               f_e\beta_v=\beta_{v'}.
                                                               \tag{136.S14}
\]
Every $\beta_v$ is an isomorphism there, since all
original arrows are inverted.

An arbitrary morphism in the ordinary localization
is represented by a finite word of arrows and
their formal inverses. For any two such words with
the same endpoints, take the finite graph consisting
of all their objects and arrows. Equation (136.S14)
rewrites either word as
$\beta_{\rm end}\beta_{\rm start}^{-1}$, so the words
are equal. Taking two isolated vertices supplies
an isomorphism between any two objects. Finally,
the category is nonempty by AQ-136-DP-MODEL.
It is therefore equivalent to the terminal groupoid.
No higher coherence assertion about an
$\infty$-categorical localization is required or
inferred: the common cone and its edge homotopies
are used to compare arbitrary \emph{finite words}
in the ordinary all-arrow groupoid. \hfill$\square$

